\documentclass[letterpaper,11pt]{article}

\usepackage[T1]{fontenc}
\usepackage[utf8]{inputenc}
\usepackage{latexsym}
\usepackage[empty]{fullpage}
\usepackage{titlesec}
\usepackage[usenames,dvipsnames]{xcolor}
\usepackage{verbatim}
\usepackage{enumitem}
\usepackage[hidelinks]{hyperref}
\usepackage[dutch,english]{babel}
\usepackage{microtype}
\usepackage{charter}

% Ensure that generated PDF is machine readable / ATS parsable
\input{glyphtounicode}
\pdfgentounicode=1

%---------- COLOR DEFINITIONS ----------
\definecolor{cvblue}{HTML}{0E5484}
\definecolor{black}{HTML}{130810}
\definecolor{darkcolor}{HTML}{0F4539}
\definecolor{airforceblue}{rgb}{0.36, 0.54, 0.66}
\definecolor{linkblue}{HTML}{1A5276}

%---------- MARGINS ----------
\addtolength{\oddsidemargin}{-0.6in}
\addtolength{\evensidemargin}{-0.6in}
\addtolength{\textwidth}{1.2in}
\addtolength{\topmargin}{-.8in}
\addtolength{\textheight}{1.6in}
\urlstyle{same}

\raggedbottom
\raggedright

%---------- SECTION FORMATTING (CLEAN, ATS-SAFE) ----------
\titleformat{\section}{
  \vspace{-5pt}\scshape\raggedright\large\bfseries\color{airforceblue}
}{}{0em}{}[\color{black}\titlerule \vspace{-4pt}]

%---------- CUSTOM ATS-OPTIMIZED COMMANDS ----------
\newcommand{\resumeSubheading}[4]{
  \vspace{-3pt}\item
    \textbf{\large#1} \hfill \textbf{\normalsize#2}\\
    \textit{\normalsize#3} \hfill \textit{\small#4}
  \vspace{-7pt}
}

\newcommand{\resumeProjectHeading}[2]{
  \vspace{-3pt}\item
    \textbf{\large#1} \hfill \textbf{\normalsize#2}\\
    \vspace{-7pt}
}

\newcommand{\resumeItem}[1]{
  \item\normalsize{#1}
}

\renewcommand\labelitemi{\textbullet}
\renewcommand\labelitemii{\textbullet}

\newcommand{\resumeSubHeadingListStart}{\begin{itemize}[leftmargin=0.0in, label={}, topsep=0pt, itemsep=2.5pt]}
\newcommand{\resumeSubHeadingListEnd}{\end{itemize}}
\newcommand{\resumeItemListStart}{\begin{itemize}[leftmargin=0.15in, topsep=0.5pt, itemsep=0.5pt, parsep=0pt]}
\newcommand{\resumeItemListEnd}{\end{itemize}}

% ==============================================================================
%                           RESUME STARTS HERE
% ==============================================================================
\begin{document}

%---------- HEADING & CONTACT INFO ----------
\newcommand{\targetRole}{Biologie \& Medisch Laboratoriumonderzoek \textbar{} Life Science \& Technology \textbar{} Moleculaire Biologie}

\begin{center}
  \vspace{-12pt}
  {\huge{Sha\"{\i}s Ayush Hemradj Ganpat}} \\[2pt]
  {\textbf{\color{airforceblue}\targetRole}} \\[3pt]
  \footnotesize{
    Delft, Zuid-Holland, Nederland \ \textbullet \
    +31 6 85141062 \ \textbullet \
    \href{mailto:shaisganpat_2000@hotmail.com}{\color{linkblue}shaisganpat\_2000@hotmail.com} \ \textbullet \
    \href{https://www.linkedin.com/in/shaisganpat/}{\color{linkblue}linkedin.com/in/shaisganpat}
  }
  \vspace{-8pt}
\end{center}

%----------- PROFIEL -----------
\section{Profiel}
\begin{itemize}[leftmargin=0in, label={}, topsep=0pt, itemsep=0pt]
  \item\normalsize{
    BML-student met een sterke wetenschappelijke basis in Life Science \& Technology (TU Delft / Universiteit Leiden) en praktijkervaring binnen de moleculaire biologie, medische microbiologie en data-analyse. Vaardig in laboratoriumtechnieken zoals PCR, primerontwerp, DNA-analyse en microbiologische kweken, gecombineerd met sterke didactische en organisatorische vaardigheden vanuit examentraining en studentenbestuur. Gemotiveerd om tijdens een afstudeerstage experimentele precisie en analytisch probleemoplossend vermogen in te zetten.
  }
\end{itemize}
\vspace{-8pt}

%----------- OPLEIDING -----------
\section{Opleiding}
\resumeSubHeadingListStart
\resumeSubheading
{Hogeschool Leiden}{Feb 2025 -- 2027 (Verwacht)}
{BSc. Biologie en Medisch Laboratoriumonderzoek (BML)}{Leiden, Nederland}
\resumeSubheading
{Delft University of Technology (TU Delft) \& Universiteit Leiden}{Sep 2019 -- Feb 2025}
{BSc. Life Science and Technology (Joint Degree)}{Delft / Leiden, Nederland}
\resumeItemListStart
\resumeItem{\textbf{Propedeuse behaald} (basis in genetica, biochemie, celbiologie en data-analyse in Python) en \textbf{Minor Lerarenopleiding Scheikunde} afgerond.}
\resumeItemListEnd
\vspace{-5pt}

\resumeSubheading
{Mr. Dr. J.C. de Miranda Lyceum}{Aug 2015 -- Jul 2018}
{VWO Voorbereidend Wetenschappelijk Onderwijs (Profiel: Natuur \& Techniek)}{Paramaribo, Suriname}
\resumeItemListStart
\resumeItem{Examencijfers: Scheikunde (9), Biologie (9), Natuurkunde (8), Engels (8), Nederlands (8), Wiskunde A/B (7), Wiskunde D (7).}
\resumeItemListEnd
\resumeSubHeadingListEnd
\vspace{-8pt}

%----------- VAARDIGHEDEN -----------
\section{Vaardigheden}
\begin{itemize}[leftmargin=0in, label={}, topsep=0pt, itemsep=1.2pt]
  \normalsize{\item{
    \textbf{Laboratorium- \& Moleculaire Technieken:}{ PCR, primerontwerp, DNA/RNA-extractie, agarose-gelelektroforese, microbiële kweken, aseptisch werken, Gram-kleuring, optische microscopie, bioveiligheid \& Good Laboratory Practice (GLP)} \\
    \textbf{Data-analyse \& Programmeren:}{ Python (Pandas, NumPy, Matplotlib, SciPy), dataverwerking, statistische modellering, wetenschappelijk rekenen, Microsoft Excel (geavanceerde data-analyse, draaitabellen)} \\
    \textbf{Didactiek \& Onderwijs:}{ Examentraining scheikunde (havo/vwo), curriculumontwikkeling, gedifferentieerd lesgeven, didactische feedback, studentmentoraat, formatief toetsen} \\
    \textbf{Bestuur \& Organisatie:}{ Studentenbeheer, faciliteitenbeheer, belangenbehartiging van bewoners, evenementenorganisatie, conflicthantering, multidisciplinair samenwerken} \\
    \textbf{Talenkennis:}{ Nederlands (moedertaal / vloeiend), Engels (vloeiend), Sranan Tongo (vloeiend), Hindi (basiskennis)}
  }}
\end{itemize}
\vspace{-8pt}

%----------- LABORATORIUMERVARING -----------
\section{Laboratorium- \& Onderzoekservaring}
\resumeSubHeadingListStart

\resumeSubheading
{Praktijkonderwijs BML -- Hogeschool Leiden}{Feb 2025 -- Heden}
{Laboratoriumonderzoeker in opleiding}{Leiden, Nederland}
\resumeItemListStart
\resumeItem{\textbf{Moleculaire Biologie \& Genetica:} Uitvoeren van practica gericht op genomische DNA-extractie, sequentie-specifiek primerontwerp en DNA-amplificatie via PCR, gevalideerd met behulp van agarose-gelelektroforese.}
\resumeItem{\textbf{Medische Microbiologie:} Kweken, isoleren en karakteriseren van bacteriën; toepassen van aseptische technieken, entmethoden en Gram-kleuring conform GLP- en bioveiligheidsrichtlijnen.}
\resumeItem{\textbf{Kwantitatieve Kwaliteitsborging:} Nauwkeurig documenteren van experimentele protocollen, interpreteren van spectrofotometrische metingen (Nanodrop) en troubleshooten van afwijkingen ter borging van reproduceerbaarheid.}
\resumeItemListEnd
\vspace{-5pt}

\resumeSubheading
{Anton de Kom Universiteit van Suriname (AdeKUS)}{Jul 2016 -- Aug 2016}
{Student-assistent Praktijkcursus Microbiologie}{Paramaribo, Suriname}
\resumeItemListStart
\resumeItem{\textbf{Laboratoriumvoorbereiding:} Geselecteerd door academische staf op basis van topprestatie in 2015; bereiden van voedingsbodems (agar), steriliseren van apparatuur (autoclaaf) en inrichten van steriele werkstations.}
\resumeItem{\textbf{Begeleiding van Cursisten:} Demonstreren van aseptische enttechnieken, isolatiestrepen en microscopische identificatie via Gram-kleuring aan een nieuwe studentengroep.}
\resumeItemListEnd
\vspace{-5pt}

\resumeSubheading
{Anton de Kom Universiteit van Suriname (AdeKUS)}{Jul 2015 -- Aug 2015}
{Cursist Intensieve Cursus Microbiologie}{Paramaribo, Suriname}
\resumeItemListStart
\resumeItem{\textbf{Praktische \& Theoretische Training:} Cursus afgerond in medische en milieumicrobiologie, steriel werken, bacterieel metabolisme en pathogenendetectie.}
\resumeItem{\textbf{Academische Erkenning:} Onderscheiden met het certificaat \textbf{``Most Motivated Student''} van het cohort wegens laboratoriumprecisie, actieve betrokkenheid en analytische werkhouding.}
\resumeItemListEnd

\resumeSubHeadingListEnd

\newpage

%----------- DIDACTISCHE & ONDERWIJSERVARING -----------
\section{Didactische \& Onderwijservaring}
\resumeSubHeadingListStart

\resumeSubheading
{Lyceo}{Apr 2023 -- Heden}
{Examentrainer Scheikunde \& Studiebegeleider}{Delft / Den Haag, Nederland}
\resumeItemListStart
\resumeItem{\textbf{Eindexamentraining:} Verzorgen van intensieve trainingen scheikunde voor examenklassen havo en vwo; complexe concepten zoals reactiekinetiek, chemisch evenwicht, thermodynamica en syntheses inzichtelijk maken.}
\resumeItem{\textbf{Individuele Begeleiding:} Didactisch analyseren van knelpunten, samenstellen van gerichte oefensets en trainen van systematische examenstrategieën met meetbare cijferverbeteringen als gevolg.}
\resumeItem{\textbf{Surveilleren \& Begeleiding:} Afnemen van gestandaardiseerde proefexamens, bewaken van examenprotocollen en bieden van vakinhoudelijke huiswerkbegeleiding bij bètavakken.}
\resumeItemListEnd
\vspace{-5pt}

\resumeSubheading
{Stanislascollege Westplantsoen}{Aug 2024 -- Jan 2025}
{Onderwijsstage Scheikunde (Stageminor Didactiek)}{Delft, Nederland}
\resumeItemListStart
\resumeItem{\textbf{Klassikaal Lesgeven:} Zelfstandig ontwerpen en verzorgen van interactieve scheikundelessen voor middelbare scholieren als onderdeel van de TU Delft Lerarenopleiding.}
\resumeItem{\textbf{Lesontwikkeling:} Ontwikkelen van didactisch lesmateriaal met interactieve demonstratieproeven en formatieve evaluaties ter versterking van scheikundige basisconcepten.}
\resumeItem{\textbf{Didactische Professionalisering:} Lespraktijk geëvalueerd via lesobservaties door werkplekbegeleiders en zelfevaluatie, gericht op klassenmanagement en differentiatie.}
\resumeItemListEnd
\vspace{-5pt}

\resumeSubheading
{Martial Arts Academy}{Jan 2017 -- Jun 2019}
{Taekwondotrainer \& Zwarte Band (1e Dan)}{Paramaribo / Delft}
\resumeItemListStart
\resumeItem{\textbf{Groepslessen \& Conditietraining:} Anderhalf jaar taekwondoles gegeven aan groepen jeugd en volwassenen met de focus op basistechnieken, sparren, conditie en discipline.}
\resumeItem{\textbf{Mentorschap:} Jonge sporters voorbereid op bandexamens met nadruk op doorzettingsvermogen, respect en sportiviteit.}
\resumeItemListEnd

\resumeSubHeadingListEnd
\vspace{-8pt}

%----------- BESTUUR & MAATSCHAPPELIJKE ERVARING -----------
\section{Bestuur \& Maatschappelijke Ervaring}
\resumeSubHeadingListStart

\resumeSubheading
{DUWO Studentenhuisvesting}{Aug 2021 -- Heden}
{Resident Assistant \& Studentbeheerder (Hotchpot / Van Hasseltlaan / BalPol4)}{Delft, Nederland}
\resumeItemListStart
\resumeItem{\textbf{Huurderscontact \& Bewonersbeheer:} Primair operationeel aanspreekpunt voor 100+ Nederlandse en internationale studenten; behandelen van praktische vragen en bevorderen van het studentenwelzijn.}
\resumeItem{\textbf{Beheer Gemeenschappelijke Ruimte:} Sinds 2024 verantwoordelijk voor het beheer van ontmoetingsruimte ``Hotchpot''; initiëren van sociale en netwerkactiviteiten om een hechte community op te bouwen.}
\resumeItem{\textbf{Veiligheid \& Inspecties:} Periodieke controles op brandveiligheid en techniek uitvoeren, gebreken inventariseren en direct schakelen met DUWO Verhuur, Schoonmaak en Vastgoedservice.}
\resumeItem{\textbf{Internationale Integratie:} Ontvangen van internationale studenten, geven van welkomstgesprekken en rondleidingen, en constructief bemiddelen bij bewonersgeschillen.}
\resumeItemListEnd
\vspace{-5pt}

\resumeSubheading
{WijWonen}{Apr 2025 -- Dec 2026}
{Bestuurslid -- Commissaris Koepel}{Delft, Nederland}
\resumeItemListStart
\resumeItem{\textbf{Belangenbehartiging:} Vertegenwoordigen van aangesloten Delftse studentenbewonersorganisaties in overleg met DUWO, gemeentelijke instanties en studentenraden.}
\resumeItem{\textbf{Beleidsafstemming:} Stimuleren van samenwerking tussen bewonerscommissies op het gebied van huurrechten, duurzaamheid en leefbaarheid.}
\resumeItemListEnd
\vspace{-5pt}

\resumeSubheading
{SUBEST Delft (Surinaamse Bèta-Studenten Delft)}{Sep 2020 -- Nov 2021}
{Bestuurslid -- Commissaris Educatie}{Delft, Nederland}
\resumeItemListStart
\resumeItem{\textbf{Educatieve Programmering:} Coördineren van educatieve evenementen, bètagerichte workshops en tentamentrainingen voor Surinamese studenten aan de TU Delft en Universiteit Leiden.}
\resumeItem{\textbf{Mentorproject:} Coördineren van de eerstejaarsbegeleiding ter bevordering van een succesvolle academische en sociale start in Nederland.}
\resumeItem{\textbf{Onderwijscontact:} Onderhouden van formele contacten tussen studentleden, academische opleidingen en externe onderwijspartners.}
\resumeItemListEnd

\resumeSubHeadingListEnd
\vspace{-8pt}

%----------- ONDERSCHEIDINGEN & CERTIFICATEN -----------
\section{Onderscheidingen, Certificaten \& Nevenactiviteiten}
\resumeItemListStart
\resumeItem{\textbf{Propedeusediploma (Life Science \& Technology):} Uitgereikt door TU Delft en Universiteit Leiden na succesvolle afronding van het volledige eerstejaars bachelortraject.}
\resumeItem{\textbf{Onderscheiding ``Most Motivated Student'':} Anton de Kom Universiteit van Suriname (AdeKUS) -- Toegekend tijdens de microbiologiecursus 2015 wegens uitstekende inzet en nauwkeurigheid.}
\resumeItem{\textbf{Zwarte Band 1e Dan in Taekwondo:} World Taekwondo (2012 -- 2019) -- Resultaat van 7+ jaar intensieve training, discipline en leiderschap.}
\resumeItemListEnd

\end{document}
